\color{unidarkgreen}{\subsection{Логика отключения выключателей}}
\color{black}

\begin{enumerate}

\item
Логика отключения В1(ЛВ) выводится накладкой \textbf{«ЛО В1(ЛВ)»}.

\item
Цепь отключения В1(ЛВ) выводится накладкой \textbf{«ЦО В1(ЛВ)»}. Время возврата импульса на отключение регулируется уставкой \textbf{«Tвозвр откл В1(ЛВ)»}. На отключение В1(ЛВ) формируется 
следующими сигналами: 

\begin{itemize}
    \item от сигнала «ЛО осн. РЗ»;
    \item от сигнала «ЛО рез. РЗ»;
    \item от сигнала «Откл. рез. защ».
\end{itemize}

\item
Запрет АПВ В1(ЛВ) выводится накладкой \textbf{«ЗАПВ В1(ЛВ)»}. Длительной импульса времени запрета ПАВ регулируется уставкой \textbf{«Тимп ЗАПВ В1(ЛВ)ВН»}. Запрет АПВ формируется следующими сигналами: 

\begin{itemize}
    \item от сигнала «ЛО осн. РЗ»;
    \item от сигнала «ЛО рез. РЗ»; 
    \item от сигнала «Откл. рез. защ».
\end{itemize}

\item
Алгоритм работы логики отключения В1(ЛВ) приведен на рисунке \ref{pic:LO_Q1_LV_VN}. Алгоритм работы логики отключения В2(ОВ) выполнен аналогично. Уставки ЛО выключателей приведены в таблицах \ref{app_set:lov1}, \ref{app_set:lov2}.

\begin{figure}[H]
\centering
\includegraphics[width=0.75\textwidth,height=0.75\textheight,keepaspectratio]{\fbpath/09.01 БСК/761.9901 ЛО В/img/LO.pdf}
\captionof{figure}{Функциональная схема логики отключения В1(ЛВ)ВН}\label{pic:LO_Q1_LV_VN}
\end{figure}

\end{enumerate} 